\section{Bellman 方程与值迭代、策略迭代}
\label{sec:16-Control-and-Reinforcement-Learning/02-Reinforcement-Learning/02}

\subsection{困境：轨迹的分支爆炸}

一个中等大小的决策问题：往前走 10 步，每步面临 3 个可选动作，每个动作后可能落脚到 5 个状态之一。穷举所有未来轨迹：

\[
(3 \times 5)^{10} \approx 5.77 \times 10^{11}.
\]

近六千亿条路径。一台每秒评估一百万条路径的计算机需要跑将近七天。而实际中的序贯决策问题极少只有 10 步——自动驾驶每秒要做几十次决策，每次面对的潜在未来是无穷的。

问题根源不是计算机慢。算力翻一万倍，步数只增加 5 层就回到原点。分支爆炸是指数型的，硬件进步追不上指数曲线。直接穷举在数学上正确，在计算上永远行不通。

\subsection{第一反应：只看眼前}

面对分支爆炸，最自然的替代就是贪心——每一步选即时奖励最高的动作，不往后看。

用两状态问题检验这条路。你站在状态 \(s\)，面前两个按钮：``取走''立即给 \(2\) 并结束。``等待''给 \(0\)，确定到达状态 \(g\)。在 \(g\) 中每步自动拿 \(1\) 并留在原地。折扣 \(\gamma = 0.9\)。

贪心选即时奖励大的：取走 \(2 >\) 等待 \(0\)，选取走。总回报 \(2\)。

不进贪心：先等待进入 \(g\)，然后无限享每步 \(1\) 的稳定收益：

\[
V^*(g) = 1 + 0.9 + 0.9^2 + 0.9^3 + \cdots = \frac{1}{1 - 0.9} = 10.
\]

从 \(s\) 选等待的总价值：\(0 + 0.9 \times 10 = 9\)。比取走的 \(2\) 多了整整 \(7\)。贪心白白扔掉了一笔巨大的未来收入。

这个反例干净地说清一件事：序贯决策中，``最好的第一步''不等于``最赚的第一步''。有时候最优第一步在即期看起来是亏的——它花钱买了一张通往高价值区域的门票。贪心策略永远不会为一个更好的未来而忍受眼前的亏损。

\subsection{撞墙：穷举炸了，贪心偏了}

穷举被指数炸掉——算不完。贪心被这个反例戳穿——看不远。两个最直接的办法都行不通。你卡在一个位置：知道要往远处看，却没有资源把远处算出来。

这就是动态规划诞生的精准动机。它不要求你先把全部未来算出来再回头比较。它只依赖一条简单到几乎可疑的恒等式：

\[
G_t = R_{t+1} + \gamma G_{t+1}.
\]

\subsection{回报自带递归}

这个恒等式是纯代数事实，不假设任何环境结构。时刻 \(t\) 的总回报正好是``这一步的即时奖励''加``折扣后的剩余回报''。如果有一个 oracle 能告诉你从下一步开始的最优剩余值，当前决策就退化成了单步贪心——贪的不再是即时奖励 \(r\)，而是 \(r + \gamma \times\) 最优未来。

策略 \(\pi\) 下取条件期望，按下一动作和下一状态展开：

\[
V^\pi(s) = \sum_a \pi(a \mid s) \sum_{s'} P(s' \mid s, a)
\bigl[ r(s, a, s') + \gamma V^\pi(s') \bigr].
\]

右端出现了 \(V^\pi\) 自己。这不是一次向前算完就结束的递推公式——它是一组联立方程。每个状态的值依赖别的状态的值，而那些状态的值又兜回来依赖它。写矩阵：

\[
V^\pi = r^\pi + \gamma P^\pi V^\pi,
\qquad
(I - \gamma P^\pi) V^\pi = r^\pi.
\]

\(\gamma < 1\) 且状态有限时，\(I - \gamma P^\pi\) 可逆，唯一解：

\[
V^\pi = (I - \gamma P^\pi)^{-1} r^\pi.
\]

这步很重要：策略固定后，决策问题退化成一个 Markov 链上的线性方程组。值函数不是某种需要凭空猜测的东西——方程钉死了它。显式求逆不是实际计算的首选，但公式确立了存在性和唯一性的数学事实。

\subsection{从期望方程到最优性方程}

最优状态值 \(V^*(s) = \sup_\pi V^\pi(s)\) 应该满足什么方程？

动态规划原理是这把钥匙：从 \(s\) 出发的最优路径，第一步选完动作、拿完奖励、落到 \(s'\) 之后，余下的路径必须仍是从 \(s'\) 出发的最优路径。如果不是——把余下部分换成更优的——整条路径立马就更优。与``最优''矛盾。

因此：

\[
V^*(s) = \max_a \sum_{s'} P(s' \mid s, a)
\bigl[ r(s, a, s') + \gamma V^*(s') \bigr].
\]

在有限动作集下，右端取到最大值的动作 \(a^*(s)\) 拼起来就给出一个最优确定性平稳策略：

\[
\pi^*(s) = a^*(s).
\]

动作值版把动作本身写进状态：

\[
Q^*(s, a) = \sum_{s'} P(s' \mid s, a)
\bigl[ r(s, a, s') + \gamma \max_{a'} Q^*(s', a') \bigr].
\]

知道 \(Q^*\) 后，最优动作直接是 \(\argmax_a Q^*(s, a)\)，不需要转移模型 \(P\)。后面转向 Q 学习的原因在这里：Q 值把决策与模型解耦。

\subsection{完整数值推演}

状态 \(g\) 只有一个动作。带入最优性方程：

\[
V^*(g) = 1 + 0.9 \, V^*(g) \quad\Rightarrow\quad V^*(g) = \frac{1}{1 - 0.9} = 10.
\]

不是跑完无穷求和，是解一个一元一次方程。

状态 \(s\) 两个动作的值：

\[
\begin{aligned}
Q^*(s, \text{取走}) &= 2 + 0 = 2, \\
Q^*(s, \text{等待}) &= 0 + 0.9 \times V^*(g) = 0 + 0.9 \times 10 = 9.
\end{aligned}
\]

\(\max\{2, 9\} = 9\)，所以 \(V^*(s) = 9\)，最优决策是等待。

关键洞察：\(V^*(g)\) 的数值 \(10\) 浓缩了从 \(g\) 出发的全部未来。当 \(s\) 的方程需要看 \(g\) 时，只需要一行：\(r + \gamma \times 10\)。动态规划的压缩效力就在这里——一个数代表了无限长轨迹的总价值。

改折扣因子。\(\gamma = 0.1\)：

\[
V^*(g) = \frac{1}{1 - 0.1} \approx 1.11,
\qquad
Q^*(s, \text{等待}) = 0 + 0.1 \times 1.11 \approx 0.11.
\]

取走值 \(2\) 重新胜出。\(\gamma\) 控制了有效视界：\(\gamma\) 小，只看眼前；\(\gamma\) 趋近 1，最远端的奖励也能传回。但 \(\gamma\) 越近 1，后面的值迭代收敛越慢——信息从末端传回首端，需要更多轮迭代。

\subsection{值迭代：反复应用最优性算子}

定义 Bellman 最优性算子

\[
(TV)(s) = \max_a \sum_{s'} P(s' \mid s, a)
\bigl[ r(s, a, s') + \gamma V(s') \bigr].
\]

值迭代：任意有界初值 \(V_0\)，反复做 \(V_{k+1} = T V_k\)。

这个迭代会不会在两个值之间来回振荡？今天的 max 指东、明天的 max 指西？

答案在压缩性。对任意有界 \(V, W\)，

\[
\| TV - TW \|_\infty \le \gamma \| V - W \|_\infty.
\]

最大值操作不放大差异，反而给差异打了一个 \(\gamma\) 的折扣。Banach 不动点定理给出三重保证：

\begin{itemize}
\tightlist
\item
  \(T\) 的唯一不动点就是 \(V^*\)；
\item
  任意有界初值都收敛到 \(V^*\)；
\item
  误差按几何级数：\(\| V_k - V^* \|_\infty \le \gamma^k \| V_0 - V^* \|_\infty\)。
\end{itemize}

以具体数字感受收敛速度。初始估值偏差最大 \(100\)。\(\gamma = 0.9\)：

\[
\begin{aligned}
\text{10 轮误差：} &\; 100 \times 0.9^{10} \approx 34.9, \\
\text{50 轮误差：} &\; 100 \times 0.9^{50} \approx 0.52, \\
\text{100 轮误差：} &\; 100 \times 0.9^{100} \approx 0.0027.
\end{aligned}
\]

\(\gamma = 0.99\)：

\[
\text{100 轮误差：} \; 100 \times 0.99^{100} \approx 36.6.
\]

一目了然：\(\gamma\) 从 \(0.9\) 提到 \(0.99\)，同样 100 轮迭代的质量从接近机器精度退化成几乎没动。折扣越接近 1，压缩越弱，远期奖励的传播需要成倍增加的迭代轮数。这不是实现问题——是信息传播速度的数学上限。

实际停止通常检查 Bellman 残差 \(\| TV_k - V_k \|_\infty\)。残差小能推出到不动点的误差有界：

\[
\| V_k - V^* \|_\infty \le \frac{\| TV_k - V_k \|_\infty}{1 - \gamma}.
\]

\(\gamma\) 近 1 时，分母 \(1 - \gamma\) 极小，即使残差看起来小，到真解的误差也可能仍然可观。单看相邻几轮某几个状态变化不大，不足以判断收敛。

\subsection{策略迭代：另一种节奏}

值迭代每轮做一次 Bellman 备份。策略迭代把工作拆得更明确：

\begin{enumerate}
\tightlist
\item
  策略评估：对当前 \(\pi\)，解 \(V^\pi = r^\pi + \gamma P^\pi V^\pi\)；
\item
  策略改进：贪心一步，
  \[
  \pi_{\text{new}}(s) \in \argmax_a \sum_{s'} P(s' \mid s, a)
  \bigl[ r(s, a, s') + \gamma V^\pi(s') \bigr].
  \]
\end{enumerate}

策略改进定理保证新策略不差于旧策略。证明直觉简洁：旧策略的价值已经算完了``之后继续走旧策略''的全部未来。如果第一步改成更高的一步前瞻动作，其余仍沿用旧策略的尾部，总回报至少不降。若新动作严格更优，整条策略严格改善。

对有限状态和动作集，确定性策略总数有限。每次改进要么严格改善、要么已经最优。有限步收敛。策略迭代在经历的策略数量通常远少于值迭代的迭代轮数——但每轮的评估成本可能很高。

完整评估昂贵的场景下，修正策略迭代：只跑几轮评估（不收敛到不动点）便改进。值迭代可视为这个谱系的极端——每轮只做一轮评估立即改进。``评估多精确才改进''是一个连续可调的参数。

\subsection{条件与边界}

动态规划给出了正确方程和收敛保证，依赖清晰的条件：

\begin{itemize}
\tightlist
\item
  模型已知：每次 Bellman 备份必须对 \(P(s' \mid s, a)\) 精确求期望；
\item
  状态和动作空间可扫描或期望可计算；
\item
  环境 Markov 且转移规律不变；
\item
  折扣 \(\gamma < 1\) 或终止型任务满足吸收条件。
\end{itemize}

状态数巨大时，逐状态扫描碰到维数灾难——这是离散版的固有问题，不是算法失误。连续状态装不进一张表；连续动作下 \(\max_a\) 本身可能就是高维非凸优化。动态规划给出正确方程，不保证方程在现实规模上可解。

去掉模型已知这一项，方程不变，但不能再精确求期望。一次样本的

\[
R_{t+1} + \gamma \max_{a'} Q(S_{t+1}, a')
\]

成为 Bellman 右端的有噪声观测。把确定性不动点迭代换成随机逼近，就得到 Q 学习。方程同一个，信息从全知变成了采样。
